\documentclass[aspectratio=169]{beamer}

% ============================================================
% PACCHETTI
% ============================================================
\usepackage[T1]{fontenc}
\usepackage[utf8]{inputenc}
\usepackage[italian]{babel}
\usepackage{lmodern}
\usepackage{tcolorbox}
\usepackage{xcolor}

% ============================================================
% TEMA E COLORI
% ============================================================
\usetheme{Madrid}

\definecolor{legaleblue}{RGB}{20,52,100}
\definecolor{legalelightblue}{RGB}{220,232,246}
\definecolor{legalegray}{RGB}{90,90,90}
\definecolor{legalered}{RGB}{150,30,30}

\setbeamercolor{palette primary}{bg=legaleblue, fg=white}
\setbeamercolor{palette secondary}{bg=legaleblue!80, fg=white}
\setbeamercolor{palette tertiary}{bg=legaleblue!60, fg=white}
\setbeamercolor{palette quaternary}{bg=legaleblue!40, fg=white}
\setbeamercolor{structure}{fg=legaleblue}
\setbeamercolor{title}{fg=white, bg=legaleblue}
\setbeamercolor{frametitle}{fg=white, bg=legaleblue}
\setbeamercolor{block title}{fg=white, bg=legaleblue}
\setbeamercolor{block body}{fg=black, bg=legalelightblue}
\setbeamercolor{alerted text}{fg=legalered}

\setbeamerfont{title}{size=\Large, series=\bfseries}
\setbeamerfont{frametitle}{size=\normalsize, series=\bfseries}

% ============================================================
% BOX NORMA (tcolorbox)
% ============================================================
\tcbuselibrary{skins, breakable}

\newtcolorbox{normabox}[1][]{%
  colback=legalelightblue,
  colframe=legaleblue,
  fonttitle=\small\bfseries,
  title={#1},
  breakable,
  left=4pt, right=4pt, top=4pt, bottom=4pt,
  boxrule=0.8pt
}

\newtcolorbox{sentenzabox}[1][]{%
  colback=gray!10,
  colframe=legalegray,
  fonttitle=\small\bfseries\color{legalegray},
  title={#1},
  breakable,
  left=4pt, right=4pt, top=4pt, bottom=4pt,
  boxrule=0.8pt
}

% ============================================================
% METADATI
% ============================================================
\title{La raccolta illecita di dati tramite \emph{form} online}
\subtitle{Profili giuridici e tecnici della raccolta di dati personali\\mediante moduli web inadeguatamente predisposti}
\author{Martino Papa \and Alessio Minati}
\institute{Università degli Studi di Firenze\\Corso di Informatica Forense}
\date{Anno Accademico 2025/2026}

% ============================================================
% DOCUMENTO
% ============================================================
\begin{document}

% ------------------------------------------------------------
% SLIDE 1 – TITOLO
% ------------------------------------------------------------
\begin{frame}
  \titlepage
\end{frame}

% ------------------------------------------------------------
% SLIDE 2 – RIFERIMENTI NORMATIVI
% ------------------------------------------------------------
\begin{frame}{Riferimenti normativi}
  \Large
  \vspace{0.5cm}
  \begin{itemize}
    \setlength\itemsep{1.2em}
    \item \textbf{Reg. (UE) 2016/679} -- GDPR
    \item \textbf{D.Lgs. 196/2003} -- Codice Privacy
    \item \textbf{D.Lgs. 101/2018} -- Adeguamento al GDPR
  \end{itemize}
\end{frame}

% ------------------------------------------------------------
% SLIDE 3 – DATI PARTICOLARI
% ------------------------------------------------------------
\begin{frame}{Dati personali e dati particolari}

  \begin{normabox}[Art.~4, n.~1 GDPR -- Definizione di dato personale]
    \small
    ``\textbf{dato personale}'': qualsiasi informazione riguardante una persona fisica identificata o identificabile.
  \end{normabox}

  \vspace{0.3cm}

  \begin{normabox}[Art.~9 GDPR -- Trattamento di categorie particolari di dati personali ]
    \small
    \textbf{Par. 1}: ``È vietato trattare dati personali che rivelino l'origine
    razziale o etnica, le opinioni politiche, le convinzioni religiose o filosofiche,
    o l'appartenenza sindacale, nonché trattare dati genetici, dati biometrici intesi
    ad identificare in modo univoco una persona fisica, dati relativi alla salute o
    alla vita sessuale o all'orientamento sessuale della persona.''\\[0.3em]
    \textbf{Par. 2}: Il divieto non si applica nei casi ivi tassativamente previsti,
    tra cui il \emph{consenso esplicito} dell'interessato (lett.~a).
  \end{normabox}

\end{frame}

% ------------------------------------------------------------
% SLIDE 4 – NORME PER LA CREAZIONE DI FORM ONLINE
% ------------------------------------------------------------
\begin{frame}{Norme per la creazione di \emph{form} online}
  \textbf{Art. 5 GDPR} -- Principi applicabili al trattamento dei dati personali

  \vspace{0.2cm}
  \begin{columns}[T]
    \begin{column}{0.48\textwidth}
      \begin{itemize}
        \item[a)] \textbf{Liceità, correttezza e trasparenza}
        \item[b)] \textbf{Limitazione della finalità}
        \item[c)] \textbf{Minimizzazione dei dati}
      \end{itemize}
    \end{column}
    \begin{column}{0.48\textwidth}
      \begin{itemize}
        \item[d)] \textbf{Esattezza}
        \item[e)] \textbf{Limitazione della conservazione}
        \item[f)] \textbf{Integrità e riservatezza}
      \end{itemize}
    \end{column}
  \end{columns}
\end{frame}

% ------------------------------------------------------------
% SLIDE 5 – OBBLIGO DI INFORMATIVA E CONSENSO
% ------------------------------------------------------------
\begin{frame}{Obbligo di informativa e consenso}

  \textbf{Art. 13 GDPR} -- Obbligo di informativa al momento della raccolta dati

  \vspace{0.3cm}
  \textbf{Art. 7 GDPR} -- Il consenso deve essere:

  \vspace{0.2cm}
  \begin{columns}[T]
    \begin{column}{0.48\textwidth}
      \begin{itemize}
        \item[\textbf{1.}] \textbf{Libero}
        \item[\textbf{2.}] \textbf{Specifico}
      \end{itemize}
    \end{column}
    \begin{column}{0.48\textwidth}
      \begin{itemize}
        \item[\textbf{3.}] \textbf{Informato}
        \item[\textbf{4.}] \textbf{Inequivocabile}
      \end{itemize}
    \end{column}
  \end{columns}
\end{frame}

% ------------------------------------------------------------
% SLIDE 6 – RESPONSABILITÀ CIVILE (ART. 82)
% ------------------------------------------------------------
\begin{frame}{Responsabilità civile del titolare -- Art. 82 GDPR}

  \begin{block}{Principio: presunzione di responsabilità}
    L'art.~82 GDPR sancisce il diritto di ogni interessato che abbia subito un danno --- materiale o immateriale --- in conseguenza di una violazione del Regolamento, di ottenere il \textbf{risarcimento} dal titolare o dal responsabile del trattamento.
  \end{block}

  \vspace{0.4cm}
  \begin{alertblock}{Onere della prova invertito}
    Il titolare è esonerato da responsabilità \textbf{soltanto se dimostra} che l'evento dannoso non gli è in alcun modo imputabile.\\[0.3em]
    \small \emph{$\Rightarrow$ Si tratta di una presunzione di responsabilità con onere della prova a carico del titolare.}
  \end{alertblock}

\end{frame}

% ------------------------------------------------------------
% SLIDE 7 – SANZIONI AMMINISTRATIVE (ART. 83)
% ------------------------------------------------------------
\begin{frame}{Sanzioni amministrative -- Art. 83 GDPR}
  \vspace{0.3cm}
  \begin{columns}[T]
    \begin{column}{0.48\textwidth}
      \begin{block}{Sanzione ordinaria (par. 4)}
        \centering
        \textbf{fino a 10 milioni €}\\
        oppure\\
        \textbf{2\% del fatturato mondiale}\\[0.2em]
        \small (es. violazione del consenso, mancata informativa)
      \end{block}
    \end{column}
    \begin{column}{0.48\textwidth}
      \begin{alertblock}{Sanzione aggravata (par. 5)}
        \centering
        \textbf{fino a 20 milioni €}\\
        oppure\\
        \textbf{4\% del fatturato mondiale}\\[0.2em]
        \small (es. violazione dei principi base, trattamento di dati particolari senza base giuridica)
      \end{alertblock}
    \end{column}
  \end{columns}

\end{frame}

% ------------------------------------------------------------
% SLIDE 8 – PROFILO PENALISTICO: ART. 167
% ------------------------------------------------------------
\begin{frame}{Profilo penalistico -- Trattamento illecito di dati}

  \begin{normabox}[Art.~167 D.Lgs. 196/2003]
    \small
    ``1. Salvo che il fatto costituisca più grave reato, chiunque, al fine di
    trarne per sé o per altri profitto ovvero al fine di recare ad altri un danno,
    procede al trattamento di dati personali in violazione di quanto disposto
    [\ldots] è punito, se dal fatto deriva nocumento, con la \textbf{reclusione da sei
    mesi a un anno e sei mesi}.\\[0.3em]
    2. Se il fatto di cui al comma 1 riguarda dati
    genetici, biometrici o relativi alla salute, alla vita sessuale, all'orientamento
    sessuale o alle opinioni, [\ldots] la pena è la \textbf{reclusione da uno a tre anni}.''
  \end{normabox}

  \vspace{0.2cm}
  Il danno e il profitto possono essere anche di natura non economica


\end{frame}

% ------------------------------------------------------------
% SLIDE 9 – PROFILO PENALISTICO: ART. 167-BIS
% ------------------------------------------------------------
\begin{frame}{Profilo penalistico -- Comunicazione e diffusione illecita su larga scala}

  \begin{normabox}[Art.~167-bis D.Lgs. 196/2003]
    \small
    ``1. Salvo che il fatto costituisca più grave reato, chiunque comunica o diffonde,
    al fine di trarne profitto per sé o altri ovvero al fine di arrecare danno, un
    archivio automatizzato o una parte sostanziale di esso contenente dati personali
    oggetto di trattamento su larga scala, in violazione degli articoli 2-ter, 2-sexies
    e 2-octies, è punito con la \textbf{reclusione da uno a sei anni}.''
  \end{normabox}

  \vspace{0.3cm}
  Si applica in caso di \textbf{data breach deliberato}: cessione o pubblicazione non autorizzata di dati raccolti tramite \emph{form}


\end{frame}

% ------------------------------------------------------------
% SLIDE 10 – PROFILO PENALISTICO: ART. 167-TER
% ------------------------------------------------------------
\begin{frame}{Profilo penalistico -- Acquisizione fraudolenta su larga scala}

  \begin{normabox}[Art.~167-ter D.Lgs. 196/2003]
    \small
    ``1. Salvo che il fatto costituisca più grave reato, chiunque, al fine di trarne
    profitto per sé o altri ovvero di arrecare danno, acquisisca con mezzi
    fraudolenti un archivio automatizzato o una parte sostanziale di esso contenente
    dati personali oggetto di trattamento su larga scala è punito con la \textbf{reclusione
    da uno a quattro anni}.''
  \end{normabox}

  \vspace{0.3cm}
  Riguarda in particolare i \emph{form} ingannevoli (\textbf{dark patterns}) i quali inducono l'utente a fornire dati senza esserne \textbf{pienamente consapevole}

\end{frame}

% ------------------------------------------------------------
% SLIDE 11 – CASO PRATICO AMMINISTRATIVO
% ------------------------------------------------------------
\begin{frame}{Caso pratico -- Profilo amministrativo}

  \begin{sentenzabox}[Garante per la protezione dei dati personali, Provv. 29 aprile 2021 -- Fonte: Garante Privacy, Newsletter n.~487]
    \small
    Un ente pubblico aveva predisposto sul proprio sito web un modulo di
    contatto (\emph{form}) che raccoglieva nome, cognome, indirizzo e-mail
    e numero di telefono degli utenti, senza che nella pagina fosse presente
    alcuna informativa adeguata ai sensi dell'art.~13 GDPR. Il Garante ha
    accertato la violazione dei principi di liceità, correttezza e trasparenza
    (art.~5, par.~1, lett.~a GDPR) e ha irrogato una sanzione pecuniaria
    amministrativa, disponendo altresì la rettifica del \emph{form}.
  \end{sentenzabox}

\end{frame}

% ------------------------------------------------------------
% SLIDE 12 – CASO PRATICO PENALE
% ------------------------------------------------------------
\begin{frame}{Caso pratico -- Profilo penale}

  \begin{sentenzabox}[Cass. pen., sez. II, 15 gennaio 2019, n.~1285]
    \small
    La Suprema Corte ha confermato la configurabilità del reato di cui all'art.~167
    del Codice Privacy in capo a un operatore che aveva sistematicamente raccolto,
    tramite un servizio web, dati personali degli utenti senza il loro consenso e li
    aveva successivamente ceduti a terzi a fini commerciali. La Corte ha precisato
    che il \emph{fine di profitto} richiesto dalla norma è integrato anche dalla sola
    prospettiva di un vantaggio economico indiretto, e che il \emph{nocumento}
    può consistere anche in un danno non patrimoniale, quale la perdita del controllo
    sui propri dati personali.
  \end{sentenzabox}

\end{frame}

% ------------------------------------------------------------
% SLIDE 13 – GRAZIE
% ------------------------------------------------------------
\begin{frame}
  \vfill
  \begin{center}
    {\Huge\bfseries\color{legaleblue} Grazie per l'attenzione}

    \vspace{0.8cm}
    \rule{0.6\linewidth}{0.4pt}

    \vspace{0.4cm}
    {\large Martino Papa \quad | \quad Alessio Minati}\\[0.3em]
    {\small Università degli Studi di Firenze\\
    Corso di Informatica Forense -- A.A. 2025/2026}

    \vspace{0.8cm}
    \rule{0.6\linewidth}{0.4pt}

    \vspace{0.4cm}
    {\small\color{legalegray} \textit{La raccolta illecita di dati tramite \emph{form} online:\\
    profili giuridici e tecnici}}
  \end{center}
  \vfill
\end{frame}

\end{document}
